\chapter{Tử Tế Nhưng Đừng Dễ Bị Dùng}
\markboth{Tử Tế Nhưng Đừng Dễ Bị Dùng}{Be Kind, But Do Not Be Easily Used}

\section{Tử Tế Không Có Nghĩa Là Cho Ai Cũng Bước Qua Em}
\begin{truyen}{Tử Tế Không Có Nghĩa Là Cho Ai Cũng Bước Qua Em}{Kindness Does Not Mean Letting Everyone Walk Over You}
\chuhoa{N}{hiều người trẻ rất sợ bị coi là khó nên gật đầu với mọi yêu cầu.}
\textit{(Many young people fear being seen as difficult, so they say yes to every request.)}

Họ tưởng đó là tử tế.
\textit{(They think that is kindness.)}

Nhưng khi người khác quen dùng em như một chỗ tiện lợi, đó không còn là giúp đỡ lành mạnh nữa.
\textit{(But when others grow used to using you as a convenient tool, that is no longer healthy help.)}

Người tử tế thật biết giúp đúng chỗ và biết dừng đúng lúc.
\textit{(A truly kind person knows how to help at the right place and stop at the right time.)}

\end{truyen}

\begin{baihoc}
Tử tế mà không có ranh giới rất dễ thành mềm yếu. Ranh giới giữ lòng tốt khỏi bị biến thành chỗ khai thác.
\textit{(Kindness without boundaries easily becomes weakness. Boundaries protect goodness from turning into something others exploit.)}
\end{baihoc}

\begin{ghinhoanh}
\textbf{Short English to remember:}

Kindness without boundaries easily becomes weakness.

Boundaries protect goodness from turning into something others exploit.

\end{ghinhoanh}

\begin{tuhocnhanh}
\textbf{Cau hoi tu hoc / Self-study questions}

1. Van de chinh la gi? \textit{What is the main problem?}

2. Bai hoc thuc te la gi? \textit{What is the practical lesson?}

3. Mot cau tieng Anh em muon nho la cau nao? \textit{Which English sentence do you want to remember?}
\end{tuhocnhanh}
\ngancach

\section{Làm Hộ Mãi Không Phải Là Nghĩa Khí}
\begin{truyen}{Làm Hộ Mãi Không Phải Là Nghĩa Khí}{Doing Everything for Others Is Not Loyalty}
\chuhoa{C}{ó người trẻ luôn làm hộ phần việc của bạn, của đồng nghiệp, của người yêu, vì sợ bị chê là không tình nghĩa.}
\textit{(Some young people keep doing the work of friends, coworkers, or partners because they fear being called disloyal.)}

Lâu dần, họ mệt nhưng vẫn cắn răng chịu.
\textit{(Over time they become exhausted but still endure silently.)}

Giúp người khác lớn không giống với việc thay họ sống phần trách nhiệm của họ.
\textit{(Helping someone grow is not the same as living their responsibility for them.)}

Nhiều mối quan hệ hỏng vì một bên cứ cho quá mức, còn bên kia quen nhận quá mức.
\textit{(Many relationships decay because one side keeps giving too much while the other gets used to taking too much.)}

\end{truyen}

\begin{baihoc}
Người trẻ cần học cách phân biệt giữa hỗ trợ và gánh thay. Nhầm hai thứ đó rất hao đời.
\textit{(Young people must learn the difference between supporting and carrying someone. Confusing the two drains life badly.)}
\end{baihoc}

\begin{ghinhoanh}
\textbf{Short English to remember:}

Young people must learn the difference between supporting and carrying someone.

Confusing the two drains life badly.

\end{ghinhoanh}

\begin{tuhocnhanh}
\textbf{Cau hoi tu hoc / Self-study questions}

1. Van de chinh la gi? \textit{What is the main problem?}

2. Bai hoc thuc te la gi? \textit{What is the practical lesson?}

3. Mot cau tieng Anh em muon nho la cau nao? \textit{Which English sentence do you want to remember?}
\end{tuhocnhanh}
\ngancach

\section{Mềm Mỏng Không Có Nghĩa Là Mờ Ranh Giới}
\begin{truyen}{Mềm Mỏng Không Có Nghĩa Là Mờ Ranh Giới}{Gentle Does Not Mean Blurry}
\chuhoa{N}{hiều người nghĩ phải thật cứng mới giữ được mình.}
\textit{(Many people think they must be hard to protect themselves.)}

Thực ra em có thể nói rất nhẹ mà vẫn rất rõ.
\textit{(In fact, you can speak very softly and still be very clear.)}

Một câu như ``Mình không làm việc đó'' thường mạnh hơn cả bài giải thích dài dòng.
\textit{(A sentence like ``I am not doing that'' is often stronger than a long explanation.)}

Ranh giới tốt không cần nhiều kịch tính.
\textit{(A good boundary does not need drama.)}

Nó chỉ cần nhất quán.
\textit{(It only needs consistency.)}

\end{truyen}

\begin{baihoc}
Người trẻ càng hiền càng nên học nói rõ. Giọng nói nhỏ không có nghĩa là quyền từ chối nhỏ.
\textit{(The gentler a young person is, the more they need to learn clear speech. A quiet voice does not mean a small right to refuse.)}
\end{baihoc}

\begin{ghinhoanh}
\textbf{Short English to remember:}

The gentler a young person is, the more they need to learn clear speech.

A quiet voice does not mean a small right to refuse.

\end{ghinhoanh}

\begin{tuhocnhanh}
\textbf{Cau hoi tu hoc / Self-study questions}

1. Van de chinh la gi? \textit{What is the main problem?}

2. Bai hoc thuc te la gi? \textit{What is the practical lesson?}

3. Mot cau tieng Anh em muon nho la cau nao? \textit{Which English sentence do you want to remember?}
\end{tuhocnhanh}
\ngancach

\section{Người Hay Thấy Có Lỗi Thường Dễ Bị Dùng}
\begin{truyen}{Người Hay Thấy Có Lỗi Thường Dễ Bị Dùng}{People Who Feel Guilty Easily Are Easy to Use}
\chuhoa{C}{ó kiểu người chỉ cần từ chối ai một lần là thấy day dứt rất lâu.}
\textit{(Some people feel guilty for a long time after refusing someone just once.)}

Những người thao túng rất nhạy với kiểu người đó.
\textit{(Manipulative people are quick to notice that kind of person.)}

Họ không cần ép mạnh.
\textit{(They do not need to force hard.)}

Họ chỉ cần gợi cảm giác em đang tệ, đang ích kỷ, hay đang vô tình.
\textit{(They only need to suggest that you are selfish, cold, or unkind.)}

Nếu không tỉnh, em sẽ cho đi vì tội lỗi chứ không còn vì tự nguyện.
\textit{(Without awareness, you begin giving out of guilt instead of choice.)}

\end{truyen}

\begin{baihoc}
Lòng tốt cần đi cùng sự tỉnh táo. Nếu không, cảm giác có lỗi sẽ bị người khác dùng như một cái dây kéo em đi.
\textit{(Kindness must be paired with awareness. Otherwise guilt becomes a rope other people use to pull you around.)}
\end{baihoc}

\begin{ghinhoanh}
\textbf{Short English to remember:}

Kindness must be paired with awareness.

Otherwise guilt becomes a rope other people use to pull you around.

\end{ghinhoanh}

\begin{tuhocnhanh}
\textbf{Cau hoi tu hoc / Self-study questions}

1. Van de chinh la gi? \textit{What is the main problem?}

2. Bai hoc thuc te la gi? \textit{What is the practical lesson?}

3. Mot cau tieng Anh em muon nho la cau nao? \textit{Which English sentence do you want to remember?}
\end{tuhocnhanh}
\ngancach

\section{Giúp Người Mà Không Bỏ Mình}
\begin{truyen}{Giúp Người Mà Không Bỏ Mình}{Help Others Without Abandoning Yourself}
\chuhoa{C}{ó những người trẻ sau một thời gian cho đi quá nhiều bỗng thấy mình cạn sạch.}
\textit{(Some young people give too much for too long and suddenly find themselves empty.)}

Lúc đó họ cay đắng và bắt đầu ghét cả sự tử tế của chính mình.
\textit{(Then they become bitter and start resenting their own kindness.)}

Vấn đề không phải lòng tốt sai.
\textit{(The problem is not that goodness is wrong.)}

Vấn đề là nó đã được dùng mà không có giới hạn.
\textit{(The problem is that it was used without limits.)}

Một lòng tốt bền là lòng tốt không đòi em phải biến mất để nó tồn tại.
\textit{(Sustainable kindness is kindness that does not require you to disappear in order to exist.)}

\end{truyen}

\begin{baihoc}
Giúp người khác là điều đẹp. Nhưng đừng làm điều đẹp theo cách phá hỏng chính mình.
\textit{(Helping others is beautiful. But do not do a beautiful thing in a way that destroys you.)}
\end{baihoc}

\begin{ghinhoanh}
\textbf{Short English to remember:}

Helping others is beautiful.

But do not do a beautiful thing in a way that destroys you.

\end{ghinhoanh}

\begin{tuhocnhanh}
\textbf{Cau hoi tu hoc / Self-study questions}

1. Van de chinh la gi? \textit{What is the main problem?}

2. Bai hoc thuc te la gi? \textit{What is the practical lesson?}

3. Mot cau tieng Anh em muon nho la cau nao? \textit{Which English sentence do you want to remember?}
\end{tuhocnhanh}
\ngancach

\section{Từ Chối Lịch Sự}
\begin{truyen}{Từ Chối Lịch Sự}{A Polite Refusal}
\chuhoa{C}{ó lúc chỉ một câu nói lịch sự rằng mình không làm được đã cứu em khỏi rất nhiều bực dọc về sau.}
\textit{(Sometimes one polite sentence saying you cannot do it saves you from much later resentment.)}

Nhiều người tưởng từ chối là cứng và làm mất lòng.
\textit{(Many think refusal is harsh and automatically offensive.)}

Nhưng một lời từ chối rõ ràng, nhẹ nhàng thường lành hơn một lời đồng ý miễn cưỡng.
\textit{(But a clear and gentle refusal is often healthier than a reluctant yes.)}

Người tôn trọng em sẽ hiểu giới hạn đó.
\textit{(People who respect you will understand that boundary.)}

\end{truyen}

\begin{baihoc}
Lịch sự không có nghĩa là để ai muốn lấy gì từ em cũng được.
\textit{(Being polite does not mean letting anyone take whatever they want from you.)}
\end{baihoc}

\begin{ghinhoanh}
\textbf{Short English to remember:}

Being polite does not mean letting anyone take whatever they want from you.

\end{ghinhoanh}

\begin{tuhocnhanh}
\textbf{Cau hoi tu hoc / Self-study questions}

1. Van de chinh la gi? \textit{What is the main problem?}

2. Bai hoc thuc te la gi? \textit{What is the practical lesson?}

3. Mot cau tieng Anh em muon nho la cau nao? \textit{Which English sentence do you want to remember?}
\end{tuhocnhanh}
\ngancach

\section{Đừng Giải Thích Quá Nhiều}
\begin{truyen}{Đừng Giải Thích Quá Nhiều}{Do Not Over-Explain}
\chuhoa{C}{ó người từ chối xong vẫn giải thích rất dài vì sợ người khác hiểu lầm.}
\textit{(Some people refuse a request and then explain for too long because they fear misunderstanding.)}

Họ tưởng càng giải thích nhiều thì càng được thông cảm.
\textit{(They think more explanation creates more understanding.)}

Nhưng giải thích quá mức đôi khi chỉ mở thêm cửa cho người khác mặc cả với ranh giới của em.
\textit{(But too much explanation can open more room for others to bargain with your boundary.)}

Một câu rõ thường mạnh hơn cả đoạn dài đầy áy náy.
\textit{(One clear sentence is often stronger than a long paragraph full of guilt.)}

\end{truyen}

\begin{baihoc}
Ranh giới tốt không cần quá nhiều lời. Nó cần sự nhất quán.
\textit{(A good boundary does not need too many words. It needs consistency.)}
\end{baihoc}

\begin{ghinhoanh}
\textbf{Short English to remember:}

A good boundary does not need too many words.

It needs consistency.

\end{ghinhoanh}

\begin{tuhocnhanh}
\textbf{Cau hoi tu hoc / Self-study questions}

1. Van de chinh la gi? \textit{What is the main problem?}

2. Bai hoc thuc te la gi? \textit{What is the practical lesson?}

3. Mot cau tieng Anh em muon nho la cau nao? \textit{Which English sentence do you want to remember?}
\end{tuhocnhanh}
\ngancach

\section{Giúp Gia Đình Có Giới Hạn}
\begin{truyen}{Giúp Gia Đình Có Giới Hạn}{Helping Family with Limits}
\chuhoa{C}{ó những người trẻ vì thương nhà nên lúc nào cũng sẵn sàng gánh thêm phần việc và phần lo của cả gia đình.}
\textit{(Some young people love their family so much that they keep taking more work and worry for the whole house.)}

Họ tưởng càng mệt vì nhà thì càng chứng tỏ mình có hiếu.
\textit{(They think the more they wear themselves out for family, the more dutiful they are.)}

Nhưng yêu thương cũng cần đi cùng phân vai và sức bền, nếu không một người sẽ kiệt mà phần còn lại không đổi.
\textit{(But love also needs structure and stamina, or one person burns out while the rest changes nothing.)}

Đỡ gia đình không có nghĩa là tan vào mọi gánh nặng của gia đình.
\textit{(Helping family does not mean dissolving into every family burden.)}

\end{truyen}

\begin{baihoc}
Hiếu thảo đẹp nhất khi đi cùng sự bền lâu, không phải sự kiệt sức.
\textit{(Filial love is most beautiful when it can last, not when it burns out.)}
\end{baihoc}

\begin{ghinhoanh}
\textbf{Short English to remember:}

Filial love is most beautiful when it can last, not when it burns out.

\end{ghinhoanh}

\begin{tuhocnhanh}
\textbf{Cau hoi tu hoc / Self-study questions}

1. Van de chinh la gi? \textit{What is the main problem?}

2. Bai hoc thuc te la gi? \textit{What is the practical lesson?}

3. Mot cau tieng Anh em muon nho la cau nao? \textit{Which English sentence do you want to remember?}
\end{tuhocnhanh}
\ngancach

\section{Tử Tế Nơi Làm Việc Không Phải Ngô Nghêch}
\begin{truyen}{Tử Tế Nơi Làm Việc Không Phải Ngô Nghêch}{Being Kind at Work Is Not Being Naive}
\chuhoa{Ở}{ chỗ làm, có người giữ thái độ dễ chịu nhưng lại học thêm cách ghi dấu công việc và nói rõ phần của mình.}
\textit{(At work, some people stay kind while also learning to record their work and speak clearly about their part.)}

Nhiều người tưởng muốn không bị dùng thì phải trở nên lạnh hoặc sắc quá mức.
\textit{(Many think the only way not to be used is to become cold or overly sharp.)}

Thật ra em vẫn có thể hiền mà rõ, tốt mà có nguyên tắc.
\textit{(In truth, you can stay gentle and still be clear, good and still principled.)}

Bản lĩnh không luôn ồn ào.
\textit{(Strength is not always loud.)}

Có khi nó chỉ nằm ở một giọng nói mềm mà ranh giới rất chắc.
\textit{(Sometimes it sits in a soft voice with a firm boundary.)}

\end{truyen}

\begin{baihoc}
Người tử tế bền là người giữ được lòng tốt mà không để lòng tốt nuốt mất mình.
\textit{(Sustainable kindness keeps its goodness without letting goodness swallow the self.)}
\end{baihoc}

\begin{ghinhoanh}
\textbf{Short English to remember:}

Sustainable kindness keeps its goodness without letting goodness swallow the self.

\end{ghinhoanh}

\begin{tuhocnhanh}
\textbf{Cau hoi tu hoc / Self-study questions}

1. Van de chinh la gi? \textit{What is the main problem?}

2. Bai hoc thuc te la gi? \textit{What is the practical lesson?}

3. Mot cau tieng Anh em muon nho la cau nao? \textit{Which English sentence do you want to remember?}
\end{tuhocnhanh}
\ngancach

\section{Nghỉ Ngơi Mà Không Có Lỗi}
\begin{truyen}{Nghỉ Ngơi Mà Không Có Lỗi}{Resting Without Guilt}
\chuhoa{C}{ó người cứ nghỉ một chút là thấy day dứt như mình đang ích kỷ hoặc lười biếng.}
\textit{(Some people feel guilty each time they rest, as if they are selfish or lazy.)}

Họ tưởng giúp người khác nghĩa là lúc nào cũng phải sẵn và hết mình.
\textit{(They think helping others means always being available and fully open.)}

Nhưng một người cạn sạch thì lòng tốt của họ cũng méo dần và chua dần.
\textit{(But a person who is empty inside will slowly distort even their kindness.)}

Nghỉ để hồi lại không phải phản bội ai.
\textit{(Resting to recover is not betrayal.)}

Đó là cách giữ cho mình còn sức mà thương đúng.
\textit{(It is how you keep enough strength to care well.)}

\end{truyen}

\begin{baihoc}
Muốn tử tế lâu, em phải cho mình quyền nghỉ.
\textit{(If you want to stay kind for long, you must allow yourself to rest.)}
\end{baihoc}

\begin{ghinhoanh}
\textbf{Short English to remember:}

If you want to stay kind for long, you must allow yourself to rest.

\end{ghinhoanh}

\begin{tuhocnhanh}
\textbf{Cau hoi tu hoc / Self-study questions}

1. Van de chinh la gi? \textit{What is the main problem?}

2. Bai hoc thuc te la gi? \textit{What is the practical lesson?}

3. Mot cau tieng Anh em muon nho la cau nao? \textit{Which English sentence do you want to remember?}
\end{tuhocnhanh}
